\documentclass[11pt, a4paper]{article}

%====================================================================
% Document Setup & Variables 
%====================================================================
\newcommand{\coursecode}{CS6886}
\newcommand{\coursename}{Systems Engineering for Deep Learning}
\newcommand{\assignmentnum}{Assignment 4: Simulating a Converted SNN}
\newcommand{\rollnum}{EE23B141}

%====================================================================
% Packages
%====================================================================
% Page Geometry Setup
\usepackage[margin=1in]{geometry}

% Typography & Fonts (Professional Research Look)
\usepackage[T1]{fontenc}
\usepackage{mathpazo} % Palatino font
\usepackage{microtype} % Improves text justification and kerning
\usepackage{setspace}
\setstretch{1.05} % Slight line spacing increase for readability

% Math & Symbols
\usepackage{amsthm, amsmath, amssymb}
\usepackage{marvosym, wasysym}

% Graphics & Tables
\usepackage{graphicx}
\usepackage{subcaption}
\usepackage{booktabs}
\usepackage{caption}
\captionsetup{labelfont=bf, font=small, skip=8pt} % Bold figure/table labels

% Code Listings 
\usepackage{listings}
\usepackage{xcolor}

% Header and Footer Control
\usepackage{fancyhdr}

% Colors
\definecolor{wordblue}{RGB}{46, 116, 181} % MS Word style slate/pale blue
\definecolor{darkwordblue}{RGB}{26, 82, 118} % Slightly darker variant for main title
\definecolor{darkblue}{rgb}{0.0, 0.15, 0.55}

% Section Headings (With Word Blue Accent)
\usepackage{titlesec}
\titleformat{\section}{\large\bfseries\color{wordblue}}{\thesection.}{1em}{}
\titleformat{\subsection}{\normalsize\bfseries\color{wordblue}}{\thesubsection.}{1em}{}

% Hyperlinks setup
\usepackage[colorlinks = true, linkcolor = darkblue, citecolor = darkblue, urlcolor = darkblue, final]{hyperref}

%====================================================================
% Custom Formatting & Macros
%====================================================================
% Professional Code block styling (Muted colors for academia)
\definecolor{codeblue}{rgb}{0.1,0.2,0.6}
\definecolor{codegreen}{rgb}{0.1,0.5,0.1}
\definecolor{codegray}{rgb}{0.5,0.5,0.5}
\definecolor{backcolour}{gray}{0.97}

\lstdefinestyle{academic}{
    backgroundcolor=\color{backcolour},   
    commentstyle=\itshape\color{codegray},
    keywordstyle=\bfseries\color{codeblue},
    numberstyle=\tiny\color{codegray},
    stringstyle=\color{codegreen},
    basicstyle=\ttfamily\footnotesize,
    breakatwhitespace=false,         
    breaklines=true,                 
    captionpos=b,                    
    keepspaces=true,                 
    numbers=left,                    
    numbersep=8pt,                  
    showspaces=false,                
    showstringspaces=false,
    showtabs=false,                  
    tabsize=4,
    frame=single,
    rulecolor=\color{lightgray}
}
\lstset{style=academic}

%====================================================================
% Header & Footer Setup
%====================================================================
\pagestyle{fancy}
\fancyhf{} % Clear default header/footer

% Make the rules aesthetic and colored
\renewcommand{\headrulewidth}{0.6pt}
\renewcommand{\footrulewidth}{0.6pt} 
% Remove the box height wrapper so the line sits naturally right under the text
\renewcommand{\headrule}{\color{wordblue}\hrule height \headrulewidth}
\renewcommand{\footrule}{\color{wordblue}\hrule height \footrulewidth}

% Word blue header text and footer page numbers
\lhead{\small\color{wordblue} \coursecode: \coursename}
\rhead{\small\color{wordblue} \rollnum}
\cfoot{\color{wordblue} \thepage}

\fancypagestyle{firstpage}{
  \fancyhf{}
  \cfoot{\color{wordblue} \thepage}
  \renewcommand{\headrulewidth}{0pt} % No top rule on the first page
  \renewcommand{\footrulewidth}{0.6pt} % Bottom rule on the first page
}

%====================================================================
% Document Body
%====================================================================
\begin{document}
\thispagestyle{firstpage}

% Title Block
\begin{center}
    \vspace*{-4.5em} 
    {\huge \bfseries \color{darkwordblue} \assignmentnum \par} % Slightly darker accent for main heading
    \vspace{1em}
    {\large Nitin Sai Maradani [\rollnum] \par} 
    \vspace{0.5em}
    {\normalsize \textit{\color{wordblue}\coursecode: \coursename} \par} % Accent added to course name
    \vspace{0.2em}
    {\normalsize \today \par}
    \vspace{0.25em}
\end{center}

% Abstract (commented out)
% \begin{abstract}
% \noindent
% This report details the methodology and results for compressing MobileNetV2 trained on the CIFAR-10 dataset. We explore the trade-offs between model quantization (weight and activation bit-widths) and global unstructured pruning (sparsity) to evaluate the impact on theoretical model size and inference accuracy. 
% \end{abstract}

\vspace{-0.5em}
{\color{wordblue}\hrule height 1pt} % Colored and slightly thicker separator
\vspace{1em} % Re-adjusted space so the first heading breathes well

% -------------------------------------------------------------------
% Main Content Begins
% -------------------------------------------------------------------

\section{Task 1: Environment Setup and Baseline Evaluation}

The assignment uses SpikingJelly, a PyTorch-based spiking neural network library, to simulate converted SNNs. After successful installation via \texttt{setup.py}, all imports executed smoothly and the pre-trained 3-layer CNN model was loaded.

The CNN architecture consists of three convolutional blocks:
\begin{itemize}
    \item \textbf{Block structure:} Each block contains \texttt{Conv2d}, \texttt{BatchNorm2d}, \texttt{ReLU}, and \texttt{AvgPool2d} modules
    \item \textbf{Input:} 28 × 28 grayscale MNIST images (1 channel)
    \item \textbf{Conv layers:} 1 → 32 → 32 → 32 channels with 3×3 kernels
    \item \textbf{Pooling:} 2×2 average pooling with stride 2 after each block
    \item \textbf{Classification head:} Flatten + Linear(32, 10)
\end{itemize}

The baseline ANN was evaluated using standard continuous FP32 inference (T = None) to establish the performance target for conversion quality:

\textbf{Baseline CNN Accuracy: 98.70\%}

This ground-truth accuracy serves as the reference for evaluating energy-accuracy trade-offs in converted SNNs.

\section{Task 2: MAC Operations Per Convolutional Layer}

To estimate computational cost, I calculated the number of Multiply-Accumulate (MAC) operations for each convolutional layer. The MAC count accounts for kernel operations over all spatial output positions:

\[
\text{MACs} = C_{\text{in}} \times k_H \times k_W \times H_{\text{out}} \times W_{\text{out}} \times C_{\text{out}}
\]

where spatial output dimensions are computed as:
\[
H_{\text{out}} = \left\lfloor \frac{H_{\text{in}} - k_H + 2p}{s} \right\rfloor + 1
\]

The automated script iterates through \texttt{model.network}, extracting kernel size, stride, padding, and channel counts from each \texttt{Conv2d} layer, while tracking spatial dimension changes through \texttt{AvgPool2d} layers.

\textbf{Results (per test image):}
\begin{table}[htbp]
  \centering
  \caption{MAC Operations Per Convolutional Layer}
  \label{tab:mac_ops}
  \renewcommand{\arraystretch}{1.2}
  \begin{tabular}{l c c c}
    \toprule
    \textbf{Layer} & \textbf{Dimensions} & \textbf{Output Size} & \textbf{MACs} \\
    \midrule
    Conv0 & 1→32, 3×3, 26×26 & 21,632 neurons & 194,688 \\
    Conv1 & 32→32, 3×3, 11×11 & 3,872 neurons & 1,115,136 \\
    Conv2 & 32→32, 3×3, 3×3 & 288 neurons & 82,944 \\
    \midrule
    \textbf{Total} & & & \textbf{1,392,768} \\
    \bottomrule
  \end{tabular}
\end{table}

At 4.6 nJ per MAC, the total ANN computational energy is \textbf{6,406,732.80 nJ} per inference.

\section{Task 3: Comparison of Conversion Scaling Modes}

ANN-to-SNN conversion replaces ReLU activations with \texttt{[VoltageScaler(1/λ) → IFNode(v\_threshold=1.0) → VoltageScaler(λ)]}, where λ is the scaling factor. The \texttt{mode} parameter determines how λ is calibrated:

\begin{itemize}
    \item \textbf{max}: λ = max(activation) across calibration batches — highest sparsity, slowest convergence
    \item \textbf{99.9\% max}: λ = 99.9th percentile (RobustNorm) — balanced sparsity and accuracy
    \item \textbf{1/2 max}: λ = 0.5 × max(activation) — moderate firing rates
    \item \textbf{1/4 max}: λ = 0.25 × max(activation) — fastest saturation, lowest sparsity
\end{itemize}

I converted the final convolutional block (split\_index = 8) using all four modes and plotted accuracy versus timesteps (T = 1 to 20).

\textbf{Observations:}
\begin{itemize}
    \item \textbf{1/4 max} reaches saturation fastest but exhibits high firing rates (poor sparsity)
    \item \textbf{max} achieves the highest sparsity but requires more timesteps to converge
    \item \textbf{99.9\% max} provides the optimal trade-off: near-baseline accuracy with good sparsity
\end{itemize}

\textbf{Design Choice:} I selected \textbf{99.9\% max (RobustNorm)} for all subsequent tasks due to its superior accuracy-sparsity balance.

\section{Task 4: Extent of Conversion}

To evaluate the impact of partial versus full conversion, I compared three scenarios:
\begin{itemize}
    \item \textbf{1 layer converted} (split\_index = 8): only Conv2 operates as SNN
    \item \textbf{2 layers converted} (split\_index = 4): Conv1 and Conv2 as SNN
    \item \textbf{3 layers converted} (split\_index = 0): full SNN conversion
\end{itemize}

The plot includes the baseline CNN accuracy (98.70\%) as a horizontal reference line.

\textbf{Results (99.9\% max, T = 20):}
\begin{table}[htbp]
  \centering
  \caption{Accuracy vs. Extent of Conversion}
  \label{tab:extent_conversion}
  \renewcommand{\arraystretch}{1.2}
  \begin{tabular}{l c c}
    \toprule
    \textbf{Configuration} & \textbf{Convergence} & \textbf{Final Accuracy} \\
    \midrule
    1 layer converted & T = 3 & 98.7\% \\
    2 layers converted & T = 5 & 98.6\% \\
    3 layers converted & T = 20 & 97.5\% \\
    \bottomrule
  \end{tabular}
\end{table}

Fully-converted SNNs exhibit an initial accuracy drop but reach saturation given sufficient timesteps, demonstrating the feasibility of complete ANN-to-SNN conversion.

\section{Task 5: Layer-wise Spike Counts and Firing Rates}

To measure spiking activity, I implemented custom PyTorch forward hooks that intercept IFNode outputs and accumulate binary spike counts across all timesteps and test images.

\textbf{Hook Design:}
\begin{itemize}
    \item Three separate hook functions (\texttt{count\_spikes\_hook0/1/2}) target each IFNode
    \item Hooks are registered before each forward pass, accumulate \texttt{output.sum().item()}, then removed
    \item Total spike counts are normalized by the test dataset size (10,000 images)
\end{itemize}

\textbf{Results (99.9\% max, T = 20, per test image):}
\begin{table}[htbp]
  \centering
  \caption{Spike Counts and Firing Rates}
  \label{tab:spike_counts}
  \renewcommand{\arraystretch}{1.2}
  \begin{tabular}{l c c c}
    \toprule
    \textbf{Layer} & \textbf{Neurons} & \textbf{Spike Count} & \textbf{Firing Rate} \\
    \midrule
    IFNode0 & 21,632 & 16,661.06 & 3.85\% \\
    IFNode1 & 3,872 & 5,493.08 & 7.09\% \\
    IFNode2 & 288 & 673.97 & 11.70\% \\
    \bottomrule
  \end{tabular}
\end{table}

\textbf{Observation:} Firing rates increase in deeper layers despite fewer neurons, reflecting higher per-neuron activity as spatial resolution decreases and feature abstraction intensifies.

\section{Task 6: Best Accuracy for ≥ 90\% Energy Savings}

Energy savings in SNNs arise from replacing dense MAC operations (4.6 nJ) with sparse event-driven Accumulate (AC) operations (0.6 nJ). The AC count for layer L+1 is:

\[
\text{ACs}_{L+1} = \text{Firing\_Rate}_L \times T \times \text{MACs}_{L+1}
\]

\textbf{Assumption:} Conv0 uses fixed MACs (4.6 nJ each) because its input is continuous FP32, not time-series spikes. Only Conv1 and Conv2 benefit from AC-based savings.

I swept T from 20 down to 15, computing:
\begin{align*}
\text{SNN Energy} &= \text{MACs}_{\text{Conv0}} \times 4.6 + (\text{ACs}_{\text{Conv1}} + \text{ACs}_{\text{Conv2}}) \times 0.6 \\
\text{Ideal Savings} &= \frac{\text{ANN Energy (Conv1+Conv2)} - \text{SNN Energy (Conv1+Conv2)}}{\text{ANN Energy (Conv1+Conv2)}}
\end{align*}

\textbf{Result:} At \textbf{T = 19}:
\begin{itemize}
    \item \textbf{Accuracy:} 97.33\%
    \item \textbf{SNN Energy:} 1,443,895.43 nJ
    \item \textbf{ANN Energy:} 6,406,732.80 nJ
    \item \textbf{Ideal Savings (Conv1+Conv2):} \textbf{90.05\%}
    \item \textbf{Total Net Savings:} 77.46\%
\end{itemize}

T = 19 is the minimum timestep configuration achieving ≥ 90\% energy savings on converted layers while maintaining high accuracy (97.33\%).

\section{Task 7: Best Energy Savings for < 5\% Accuracy Loss}

The constraint is: SNN accuracy ≥ 93.70\% (baseline 98.70\% minus 5\%).

I swept T from 15 down to 11:

\begin{table}[htbp]
  \centering
  \caption{Accuracy and Energy Savings vs. Timesteps}
  \label{tab:energy_sweep}
  \renewcommand{\arraystretch}{1.2}
  \begin{tabular}{c c c c c}
    \toprule
    \textbf{T} & \textbf{Accuracy} & \textbf{Loss (\%)} & \textbf{Ideal Savings} & \textbf{Constraint} \\
    \midrule
    15 & 95.47\% & 3.23\% & 92.58\% & ✓ \\
    \textbf{14} & \textbf{94.71\%} & \textbf{3.99\%} & \textbf{93.26\%} & ✓ \\
    13 & 92.48\% & 6.22\% & 93.95\% & ✗ \\
    \bottomrule
  \end{tabular}
\end{table}

\textbf{Result:} At \textbf{T = 14}:
\begin{itemize}
    \item \textbf{Accuracy:} 94.71\% (loss = 3.99\% < 5\%)
    \item \textbf{Ideal Savings (Conv1+Conv2):} \textbf{93.26\%}
    \item \textbf{Total Net Savings:} 80.22\%
\end{itemize}

T = 14 maximizes energy savings while satisfying the accuracy degradation bound.

\section{Task 8: Memory-Aware CNN Write Costs}

Real hardware must write intermediate activation maps to memory. Assuming high-bandwidth memory with 25 nJ per 64-byte write, I computed the cost for each fused block output.

\textbf{Fused Block Output Sizes (after AvgPool):}
\begin{itemize}
    \item Fused0: 13×13×32 = 5,408 activations
    \item Fused1: 5×5×32 = 800 activations
    \item Fused2: 1×1×32 = 32 activations
\end{itemize}

For FP32 (4 bytes/activation), 16 activations fit per 64-byte packet:
\[
\text{Write Cost} = \left\lceil \frac{\text{Activation Count}}{16} \right\rceil \times 25 \text{ nJ}
\]

\begin{table}[htbp]
  \centering
  \caption{CNN Memory Write Costs}
  \label{tab:cnn_memory}
  \renewcommand{\arraystretch}{1.2}
  \begin{tabular}{l c c c}
    \toprule
    \textbf{Block} & \textbf{Activations} & \textbf{Packets} & \textbf{Write Cost (nJ)} \\
    \midrule
    Fused0 & 5,408 & 338 & 8,450 \\
    Fused1 & 800 & 50 & 1,250 \\
    Fused2 & 32 & 2 & 50 \\
    \midrule
    \textbf{Total} & & & \textbf{9,750} \\
    \bottomrule
  \end{tabular}
\end{table}

\section{Task 9: Memory-Aware SNN Energy Model}

SNNs emit binary spikes (1 bit/activation), packing 512 activations per 64-byte packet:
\[
\text{SNN Write Cost per timestep} = \left\lceil \frac{\text{Activation Count}}{512} \right\rceil \times 25 \text{ nJ}
\]

\begin{table}[htbp]
  \centering
  \caption{SNN Memory Write Costs (per timestep)}
  \label{tab:snn_memory}
  \renewcommand{\arraystretch}{1.2}
  \begin{tabular}{l c c c}
    \toprule
    \textbf{Block} & \textbf{Activations} & \textbf{Packets} & \textbf{Write Cost (nJ)} \\
    \midrule
    Fused0 & 5,408 & 11 & 275 \\
    Fused1 & 800 & 2 & 50 \\
    Fused2 & 32 & 1 & 25 \\
    \midrule
    \textbf{Total} & & & \textbf{350} \\
    \bottomrule
  \end{tabular}
\end{table}

At T = 14 (the optimal operating point from Task 7):

\begin{align*}
\text{CNN Total Energy} &= 6,406,732.80 + 9,750 = 6,416,482.80 \text{ nJ} \\
\text{SNN Total Energy} &= 1,267,093.59 + (350 \times 14) = 1,267,093.59 + 4,900 = 1,271,993.59 \text{ nJ}\\
\text{Net Savings} &= \frac{6,416,482.80 - 1,271,993.59}{6,416,482.80} = \textbf{80.18\%}
\end{align*}

\textbf{Ideal Savings (Conv1+Conv2 only):} \textbf{93.42\%}

\textbf{Conclusion:} Even with memory overhead, the 1-bit spike representation maintains substantial energy savings. This highlights why neuromorphic innovations like \textbf{in-memory compute} and \textbf{event-driven asynchronous architectures} (which eliminate off-chip activation writes entirely) are critical to fully realizing SNN efficiency gains.

%====================================================================
% References (Optional but recommended)
%====================================================================
\vspace{2em}
\begin{thebibliography}{9}
\bibitem{spikingjelly}
Fang, W., Chen, Y., Ding, J., Yu, Z., Masquelier, T., Chen, D., Huang, L., Zhou, H., Li, G., \& Tian, Y. (2023).
\textit{SpikingJelly: An open-source machine learning platform based on PyTorch for spiking neural networks.}
Science Advances, 9(40), eadi1480.

\bibitem{rueckauer2017}
Rueckauer, B., Lungu, I. A., Hu, Y., Pfeiffer, M., \& Liu, S. C. (2017).
\textit{Conversion of continuous-valued deep networks to efficient event-driven networks for image classification.}
Frontiers in Neuroscience, 11, 682.
\end{thebibliography}

\end{document}